% ============================================================
% WarpX Polywell Simulation — Full Technical Documentation
% Compile with: pdflatex DOCUMENTATION.tex  (run twice for TOC)
% Or:           latexmk -pdf DOCUMENTATION.tex
% ============================================================
\documentclass[11pt, a4paper]{article}

% --- Packages ---
\usepackage[utf8]{inputenc}
\usepackage[T1]{fontenc}
\usepackage{lmodern}
\usepackage[margin=2.5cm]{geometry}
\usepackage{hyperref}
\usepackage{xcolor}
\usepackage{listings}
\usepackage{booktabs}
\usepackage{array}
\usepackage{longtable}
\usepackage{amsmath}
\usepackage{graphicx}
\usepackage{titlesec}
\usepackage{parskip}
\usepackage{fancyhdr}
\usepackage{microtype}
\usepackage{enumitem}

% --- Hyperref colours ---
\hypersetup{
    colorlinks=true,
    linkcolor=blue!70!black,
    urlcolor=blue!70!black,
    citecolor=green!50!black,
    pdftitle={WarpX Polywell Simulation Documentation},
    pdfauthor={},
}

% --- Code listing style ---
\definecolor{codebg}{gray}{0.96}
\definecolor{codeframe}{gray}{0.75}
\definecolor{kwcolor}{RGB}{0,0,180}
\definecolor{strcolor}{RGB}{163,21,21}
\definecolor{cmtcolor}{RGB}{0,128,0}

\lstset{
    backgroundcolor=\color{codebg},
    frame=single,
    rulecolor=\color{codeframe},
    basicstyle=\ttfamily\small,
    keywordstyle=\color{kwcolor}\bfseries,
    stringstyle=\color{strcolor},
    commentstyle=\color{cmtcolor}\itshape,
    breaklines=true,
    breakatwhitespace=true,
    showstringspaces=false,
    tabsize=4,
    captionpos=b,
    numbers=left,
    numberstyle=\tiny\color{gray},
    xleftmargin=1.5em,
    framexleftmargin=1.5em,
}

\lstdefinestyle{bash}{language=bash, morekeywords={conda,mpirun,python}}
\lstdefinestyle{python}{language=python}
\lstdefinestyle{plain}{language={}, numbers=none}

% --- Section title formatting ---
\titleformat{\section}{\large\bfseries\sffamily}{{\thesection}}{0.8em}{}[\vspace{-0.3em}\rule{\linewidth}{0.4pt}]
\titleformat{\subsection}{\normalsize\bfseries\sffamily}{\thesubsection}{0.8em}{}
\titleformat{\subsubsection}{\normalsize\bfseries}{\thesubsubsection}{0.8em}{}

% --- Header / footer ---
\pagestyle{fancy}
\fancyhf{}
\lhead{\sffamily\small WarpX Polywell Simulation}
\rhead{\sffamily\small Documentation}
\cfoot{\thepage}
\renewcommand{\headrulewidth}{0.4pt}

% --- Table helpers ---
\newcolumntype{L}[1]{>{\raggedright\arraybackslash}p{#1}}

% ============================================================
\begin{document}

\begin{titlepage}
    \centering
    \vspace*{3cm}
    {\Huge\bfseries\sffamily WarpX Polywell Simulation\\[0.4em]}
    {\Large\sffamily Technical Documentation\\[2em]}
    \rule{0.6\linewidth}{0.6pt}\\[1.5em]
    {\large Setup, Architecture, and Coding Reference}\\[3em]
    {\normalsize Last updated: 2026-03-18}
    \vfill
\end{titlepage}

\tableofcontents
\newpage

% ============================================================
\section{Project Overview}

This project uses the \href{https://ecp-warpx.github.io/}{WarpX} particle-in-cell (PIC)
framework to simulate plasma confinement in a \textbf{polywell} magnetic configuration.

\subsection*{Workflow}
\begin{enumerate}
    \item \textbf{Generate external field grids} (B and E) using analytical/numerical
          methods and write them to openPMD-compliant HDF5 files.
    \item \textbf{Run a PIC simulation} — WarpX reads those grids at startup, then
          self-consistently evolves particle trajectories and electromagnetic fields.
    \item \textbf{Output diagnostics} — field snapshots and particle data are written
          every 100 steps in openPMD/HDF5 format.
\end{enumerate}

\subsection*{Key Design Decisions}

\begin{longtable}{L{4.5cm} L{9.5cm}}
\toprule
\textbf{Decision} & \textbf{Rationale} \\
\midrule
Field pre-computation + caching &
    B and E fields are expensive to compute. Results are saved to \texttt{.h5} files;
    the filename encodes all parameters as a cache key so repeated runs skip computation. \\[0.4em]
openPMD 1.1.0 format &
    Ensures compatibility with WarpX's \texttt{read\_from\_file} interface,
    \texttt{openpmd\_viewer}, VisIt, and ParaView. \\[0.4em]
Non-vectorised E-field loop &
    The analytic E-field integrand has per-coil coordinate transforms that are
    difficult to batch; each of the $N^3 \times 6$ evaluations runs in a plain Python loop. \\[0.4em]
\texttt{do\_initial\_div\_cleaning = 0} &
    The magpylib B-field is analytically divergence-free ($\nabla \cdot \mathbf{B} = 0$).
    Divergence cleaning with open boundary conditions introduces numerical errors
    at startup and is unnecessary. \\
\bottomrule
\end{longtable}

\subsection*{Quick Start}
\begin{lstlisting}[style=bash]
conda activate warpx-env
cd /path/to/WarpX
python inputs/polywell_input.py
\end{lstlisting}

% ============================================================
\section{Installation \& Environment Setup}

\subsection{Prerequisites}
\begin{itemize}
    \item \href{https://docs.conda.io/en/latest/miniconda.html}{Miniconda} or Anaconda
    \item Python 3.9+ (managed through conda)
    \item 8+ CPU cores recommended (\texttt{N} must be divisible by the MPI rank count)
\end{itemize}

\subsection{Step 1 — Create a Conda Environment}
\begin{lstlisting}[style=bash]
conda create -n warpx-env python=3.10
conda activate warpx-env
\end{lstlisting}

\subsection{Step 2 — Install WarpX via conda-forge}
\begin{lstlisting}[style=bash]
conda install -c conda-forge warpx
\end{lstlisting}
Installs \texttt{pywarpx} (the PICMI Python interface) and all compiled WarpX
dependencies (MPI, HDF5, ADIOS2, etc.).

\subsection{Step 3 — Install Additional Packages}
\begin{lstlisting}[style=bash]
conda install -c conda-forge magpylib h5py
\end{lstlisting}

\begin{longtable}{L{3cm} L{10cm}}
\toprule
\textbf{Package} & \textbf{Purpose} \\
\midrule
\texttt{magpylib} & Magnetic field calculation for current loop coils \\
\texttt{h5py}     & Read/write HDF5 field grid files \\
\texttt{numpy}    & Array math (installed as a WarpX transitive dependency) \\
\bottomrule
\end{longtable}

\subsection{Step 4 — Verify}
\begin{lstlisting}[style=python]
import pywarpx
from pywarpx import picmi
import magpylib
import h5py
print("All imports OK")
\end{lstlisting}

\subsection{Troubleshooting}

\begin{longtable}{L{5.5cm} L{8cm}}
\toprule
\textbf{Symptom} & \textbf{Fix} \\
\midrule
\texttt{ModuleNotFoundError: src} &
    Run from project root: \texttt{cd /path/to/WarpX} \\[0.3em]
MPI domain decomposition error &
    \texttt{N} must be divisible by rank count. Default \texttt{N=72} divides by
    1, 2, 3, 4, 6, 8, 9, 12\ldots \\[0.3em]
Divergence cleaning crash at startup &
    Keep \texttt{warpx.do\_initial\_div\_cleaning = 0} \\
\bottomrule
\end{longtable}

% ============================================================
\section{Project Structure}

\begin{lstlisting}[style=plain, caption={Directory tree}]
WarpX/
|-- docs/                        # Documentation (this directory)
|   |-- DOCUMENTATION.md         # Combined Markdown reference
|   |-- DOCUMENTATION.tex        # This LaTeX source
|   |-- README.md                # Navigation index
|   |-- packages                 # Conda packages used during development
|   |-- setup/
|   |   |-- installation.md
|   |   `-- project-structure.md
|   |-- simulation/
|   |   |-- running.md
|   |   |-- parameters.md
|   |   `-- output.md
|   |-- modules/
|   |   |-- bext.md
|   |   |-- eext.md
|   |   `-- utils.md
|   `-- physics/
|       `-- polywell.md
|
|-- inputs/
|   `-- polywell_input.py        # Main entry point (edit parameters here)
|
|-- output/
|   `-- bext/                    # Auto-generated HDF5 field files
|       `-- *.h5
|
|-- src/
|   |-- bext/
|   |   |-- bext.py             # B-field HDF5 generation
|   |   `-- make_collection.py  # magpylib coil geometry
|   |-- eext/
|   |   |-- eext.py             # E-field HDF5 population
|   |   `-- methods.py          # Analytic E-field methods + EMethods enum
|   |-- post/
|   |   `-- reader.py           # Post-processing placeholder (empty)
|   `-- utils/
|       |-- cyl.py              # Cartesian <-> Cylindrical conversion
|       `-- paths.py            # ROOT_DIR and BEXT_DIR constants
|
`-- tests/                       # Jupyter notebooks for exploration
    |-- input_fields.ipynb
    |-- magpylib.ipynb
    `-- openPMD-viewer.ipynb
\end{lstlisting}

\subsection*{Key File Roles}

\begin{longtable}{L{5.5cm} L{8.5cm}}
\toprule
\textbf{File} & \textbf{Role} \\
\midrule
\texttt{inputs/polywell\_input.py} &
    Entry point. All user parameters. Orchestrates field file generation,
    WarpX setup, and \texttt{sim.step()}. \\[0.3em]
\texttt{src/bext/bext.py} &
    Computes 3D B-field via magpylib; writes openPMD HDF5. \\[0.3em]
\texttt{src/bext/make\_collection.py} &
    Builds 6-coil polywell \texttt{magpylib.Collection}. \\[0.3em]
\texttt{src/eext/eext.py} &
    Overwrites zeroed E-field datasets with analytically computed values. \\[0.3em]
\texttt{src/eext/methods.py} &
    \texttt{fw\_e}, \texttt{bob\_e} functions and \texttt{EMethods} enum. \\[0.3em]
\texttt{src/utils/paths.py} &
    \texttt{ROOT\_DIR} and \texttt{BEXT\_DIR}, resolved via \texttt{\_\_file\_\_}. \\[0.3em]
\texttt{src/utils/cyl.py} &
    \texttt{toCyl()} and \texttt{toCart()} helpers. \\[0.3em]
\texttt{output/bext/} &
    Cached HDF5 field files. Filename encodes all parameters as cache key. \\
\bottomrule
\end{longtable}

% ============================================================
\section{Running a Simulation}

\subsection{Basic Run}
\begin{lstlisting}[style=bash]
cd /path/to/WarpX
conda activate warpx-env
python inputs/polywell_input.py
\end{lstlisting}

\subsection{MPI Parallel Run}
\begin{lstlisting}[style=bash]
mpirun -n 8 python inputs/polywell_input.py
\end{lstlisting}
\textbf{Note:} \texttt{N} must be divisible by the rank count. Default \texttt{N=72}
works with 8 ranks.

\subsection{Execution Flow}

\begin{lstlisting}[style=plain]
polywell_input.py
|
|-- 1. Read user parameters (top of file)
|
|-- 2. Build WarpX objects
|   |-- picmi.Cartesian3DGrid       -- domain [-L,L]^3, open/absorbing BCs
|   |-- picmi.ElectromagneticSolver -- Yee scheme, CFL=0.99
|   |-- picmi.UniformDistribution   -- plasma in +/-(0.11*L)^3
|   |-- picmi.Species (electron)
|   `-- picmi.Species (proton)
|
|-- 3. make_bext_file(I, b_dia, b_offset, L, N)
|       |-- cached? -> return path
|       `-- compute with magpylib, write HDF5, return path
|
|-- 4. fill_eext_file(path, method, e_dia, e_offset, Q, L, N)  [optional]
|       |-- cached? -> return path
|       `-- compute analytically, update HDF5, rename, return path
|
|-- 5. warpx.read_fields_from_path = ext_path
|
|-- 6. picmi.Simulation + add_species + add_diagnostic
|
`-- 7. sim.step()   -- runs for max_steps
\end{lstlisting}

\subsection{Field File Caching}

\begin{longtable}{L{2.5cm} L{11.5cm}}
\toprule
\textbf{Type} & \textbf{Filename pattern} \\
\midrule
B only &
\texttt{B\_ext\_I-\{I\}A\_D-\{dia\}m\_Off-\{offset\}m\_L-\{L\}m\_N-\{N\}.h5} \\[0.3em]
B + E &
\texttt{B\_ext\_...\_ E\_ext\_Q-\{Q\}\_D-\{e\_dia\}m\_offset-\{e\_offset\}m\_C\_L-\{L\}m\_N-\{N\}.h5} \\
\bottomrule
\end{longtable}

Old files are not deleted automatically --- clean \texttt{output/bext/} manually when needed.

% ============================================================
\section{Input Parameters Reference}

All parameters live in the \texttt{=== USER INPUTS ===} block at the top of
\texttt{inputs/polywell\_input.py}.

\subsection{Simulation Control}
\begin{longtable}{L{3.5cm} L{2cm} L{1.5cm} L{6.5cm}}
\toprule
\textbf{Parameter} & \textbf{Default} & \textbf{Unit} & \textbf{Description} \\
\midrule
\texttt{max\_steps} & \texttt{10000} & steps & Total PIC time steps \\
\texttt{warpx.const\_dt} & \texttt{1e-12} & s & Fixed time step size \\
\texttt{p\_density} & \texttt{1e18} & m$^{-3}$ & Number density of each plasma species \\
\bottomrule
\end{longtable}

\subsection{B-Field — Polywell Coils}
\begin{longtable}{L{3.5cm} L{2cm} L{1.5cm} L{6.5cm}}
\toprule
\textbf{Parameter} & \textbf{Default} & \textbf{Unit} & \textbf{Description} \\
\midrule
\texttt{I}        & \texttt{1e6} & A & Current through each coil \\
\texttt{b\_dia}   & \texttt{1}   & m & Coil diameter \\
\texttt{b\_offset}& \texttt{1.1} & m & Distance from origin to coil center \\
\bottomrule
\end{longtable}

\subsection{E-Field — Charged Rings}
\begin{longtable}{L{3.5cm} L{2cm} L{1.5cm} L{6.5cm}}
\toprule
\textbf{Parameter} & \textbf{Default} & \textbf{Unit} & \textbf{Description} \\
\midrule
\texttt{e\_method} & \texttt{"FW"} & -- & Method name, or \texttt{None} to disable E-field \\
\texttt{Q}         & \texttt{1e-9} & C  & Total charge per ring \\
\texttt{e\_dia}    & \texttt{0.75} & m  & Ring diameter \\
\texttt{e\_offset} & \texttt{1.1}  & m  & Distance from origin to ring center \\
\bottomrule
\end{longtable}

\subsubsection*{Available E-field methods}
\begin{longtable}{L{2.5cm} L{3cm} L{7.5cm}}
\toprule
\textbf{Name} & \textbf{Function} & \textbf{Notes} \\
\midrule
\texttt{"FW"}  & \texttt{fw\_e()}  & 500-point Riemann sum over $\theta \in [0, 2\pi]$ \\
\texttt{"BOB"} & \texttt{bob\_e()} & 100-point normalised integration \\
\texttt{None}  & ---               & E-field disabled; zero arrays used \\
\bottomrule
\end{longtable}

\subsection{Grid}
\begin{longtable}{L{4.5cm} L{2.5cm} L{1.5cm} L{5cm}}
\toprule
\textbf{Parameter} & \textbf{Default} & \textbf{Unit} & \textbf{Description} \\
\midrule
\texttt{L} & \texttt{3} & m & Domain half-length; full domain $[-L, L]^3$ \\
\texttt{N} & \texttt{72} & -- & Grid cells per axis \\
\texttt{number\_per\_cell\_each\_dim} & \texttt{[10,10,10]} & -- & Macro-particles per cell per dimension \\
\texttt{plasma\_bounding} & \texttt{0.11} & fraction & Plasma initialised within $\pm(L \times 0.11)$ \\
\bottomrule
\end{longtable}

\subsection{Solver Options}
\begin{lstlisting}[style=python]
# Active:
solver = picmi.ElectromagneticSolver(grid=grid, method="Yee", cfl=0.99)

# Alternatives (commented out in file):
# solver = picmi.ElectrostaticSolver(grid=grid)
# solver = picmi.HybridPICSolver(...)  # does NOT support E-field from file
\end{lstlisting}

\subsection{Diagnostics}

Both write every 100 steps in openPMD/HDF5 format.

\begin{longtable}{L{3cm} L{10cm}}
\toprule
\textbf{Diagnostic} & \textbf{Data} \\
\midrule
\texttt{field\_diag} & \texttt{Bx}, \texttt{By}, \texttt{Bz}, \texttt{Jx}, \texttt{Jy}, \texttt{Jz}, \texttt{part\_per\_cell} \\
\texttt{part\_diag}  & \texttt{x}, \texttt{y}, \texttt{z}, \texttt{ux}, \texttt{uy}, \texttt{uz}, \texttt{weighting} (for \texttt{plasma\_e} and \texttt{plasma\_i}) \\
\bottomrule
\end{longtable}

% ============================================================
\section{Output Files}

\subsection{External Field Files (\texttt{output/bext/*.h5})}

Pre-computed inputs to WarpX. Generated once, cached, reused.

\subsubsection*{HDF5 Internal Layout}
\begin{lstlisting}[style=plain]
/ (root)
    attrs: openPMD="1.1.0", basePath="/data/%T/", meshesPath="meshes/", ...
`-- data/
    `-- 1/
        attrs: time=0.0, dt=0.0, timeUnitSI=1.0
        `-- meshes/
            |-- B/
            |   attrs: geometry="cartesian",
            |          gridSpacing=[dx,dy,dz],
            |          gridGlobalOffset=[-L,-L,-L],
            |          unitDimension=[0,1,1,-2,0,0,-1]
            |   |-- x  float64 (N,N,N)  unitSI=1.0 (Tesla)
            |   |-- y  float64 (N,N,N)
            |   `-- z  float64 (N,N,N)
            `-- E/
                attrs: geometry="cartesian",
                       gridSpacing=[dx,dy,dz],
                       gridGlobalOffset=[-L,-L,-L],
                       unitDimension=[1,1,-3,-1,0,0,0]
                |-- x  float64 (N,N,N)  unitSI=1.0 (V/m)
                |-- y  float64 (N,N,N)
                `-- z  float64 (N,N,N)
\end{lstlisting}

\subsubsection*{Reading with h5py}
\begin{lstlisting}[style=python]
import h5py
with h5py.File("output/bext/B_ext_....h5", "r") as f:
    Bx = f["data/1/meshes/B/x"][:]    # shape (N, N, N), Tesla
    Ex = f["data/1/meshes/E/x"][:]    # shape (N, N, N), V/m
    spacing = f["data/1/meshes/B"].attrs["gridSpacing"]
    offset  = f["data/1/meshes/B"].attrs["gridGlobalOffset"]
\end{lstlisting}

\subsubsection*{Reading with openPMD-viewer}
\begin{lstlisting}[style=python]
from openpmd_viewer import OpenPMDTimeSeries
ts = OpenPMDTimeSeries("output/bext/")
Bz, info = ts.get_field(field="B", coord="z", iteration=1)
\end{lstlisting}

\subsection{Simulation Diagnostic Files}

Written every 100 steps by WarpX during \texttt{sim.step()}.

\begin{longtable}{L{3cm} L{3.5cm} L{7cm}}
\toprule
\textbf{Diagnostic} & \textbf{Quantity} & \textbf{Description} \\
\midrule
\texttt{field\_diag} & \texttt{Bx, By, Bz} & Magnetic field (T) \\
                     & \texttt{Jx, Jy, Jz} & Current density (A/m$^2$) \\
                     & \texttt{part\_per\_cell} & Macro-particles per cell \\
\midrule
\texttt{part\_diag} & \texttt{x, y, z} & Positions (m) \\
                    & \texttt{ux, uy, uz} & Normalised momenta $\gamma v/c$ \\
                    & \texttt{weighting} & Macro-particle statistical weight \\
\bottomrule
\end{longtable}

% ============================================================
\section{Module: \texttt{src/bext}}

Generates the 3D B-field grid and writes it as an openPMD HDF5 file.

\subsection{\texttt{make\_collection.py}}

\subsubsection*{\texttt{make\_polywell\_collection(a, dia, d)} $\to$ \texttt{magpylib.Collection}}

Six circular current loops at $\pm x$, $\pm y$, $\pm z$ with alternating polarities.

\begin{longtable}{L{1.5cm} L{3cm} L{2.5cm} L{6cm}}
\toprule
\textbf{Coil} & \textbf{Position} & \textbf{Current} & \textbf{Rotation} \\
\midrule
s1 & $(-d,\,0,\,0)$ & $+a$ & 90° around Y \\
s2 & $(+d,\,0,\,0)$ & $-a$ & 90° around Y \\
s3 & $(0,\,-d,\,0)$ & $-a$ & $-90°$ around X \\
s4 & $(0,\,+d,\,0)$ & $+a$ & $-90°$ around X \\
s5 & $(0,\,0,\,-d)$ & $+a$ & 90° around Z \\
s6 & $(0,\,0,\,+d)$ & $-a$ & 90° around Z \\
\bottomrule
\end{longtable}

\subsection{\texttt{bext.py}}

\subsubsection*{\texttt{get\_bext\_file\_name(I, dia, offset, L, N)} $\to$ \texttt{str}}

Returns the expected filename for a given parameter set.
\begin{lstlisting}[style=python]
get_bext_file_name(1e6, 1.0, 1.1, 3, 72)
# -> "B_ext_I-1000000.0A_D-1m_Off-1.1m_L-3m_N-72.h5"
\end{lstlisting}

\subsubsection*{\texttt{make\_bext\_file(I, dia, offset, L, N)} $\to$ \texttt{Path}}

Main public function. Returns cached path if file exists; otherwise:
\begin{enumerate}
    \item \texttt{make\_polywell\_collection(I, dia, offset)} $\to$ \texttt{magpylib.Collection}
    \item \texttt{np.meshgrid(linspace($-L$,$L$,$N$), indexing=\textquotesingle ij\textquotesingle)} $\to$ mesh shape $(N,N,N,3)$
    \item \texttt{collection.getB(mesh)} $\to$ $B_{x}$, $B_{y}$, $B_{z}$ each $(N,N,N)$
    \item \texttt{\_make\_empty\_ext\_h5(path)} $\to$ skeleton HDF5 with openPMD metadata
    \item \texttt{\_fill\_h5\_file(...)} $\to$ writes B datasets; initialises E to zero
\end{enumerate}

\subsubsection*{openPMD attributes written by \texttt{\_make\_empty\_ext\_h5}}

\begin{longtable}{L{4.5cm} L{9cm}}
\toprule
\textbf{Attribute} & \textbf{Value} \\
\midrule
\texttt{openPMD}           & \texttt{"1.1.0"} \\
\texttt{basePath}          & \texttt{"/data/\%T/"} \\
\texttt{iterationEncoding} & \texttt{"fileBased"} \\
B \texttt{unitDimension}   & \texttt{[0, 1, 1, -2, 0, 0, -1]} (Tesla) \\
E \texttt{unitDimension}   & \texttt{[1, 1, -3, -1, 0, 0, 0]} (V/m) \\
\bottomrule
\end{longtable}

% ============================================================
\section{Module: \texttt{src/eext}}

Computes the E-field from charged rings and overwrites the zeroed placeholder
in the B-field HDF5.

\subsection{\texttt{methods.py}}

\subsubsection*{\texttt{fw\_e(r, z, a, Q, resolution=500)} $\to$ \texttt{(E\_r, E\_z)}}

Numerical integration over $\theta \in [0,\,2\pi]$ with linear charge density
$\lambda = Q/(2\pi a)$:

\begin{align}
E_r &= \frac{\lambda a}{4\pi\varepsilon_0}
       \int_0^{2\pi} \frac{r - a\cos\theta}{D^3}\,d\theta \\[6pt]
E_z &= \frac{\lambda a\,z}{4\pi\varepsilon_0}
       \int_0^{2\pi} \frac{1}{D^3}\,d\theta \\[4pt]
D   &= \sqrt{r^2 + a^2 - 2ar\cos\theta + z^2}
\end{align}

\subsubsection*{\texttt{bob\_e(r, z, a, Q, resolution=100)} $\to$ \texttt{(E\_r, E\_z)}}

Alternative normalised integration over $\theta \in [0,\,\pi]$ using
$k = 8.99\times10^9\,\text{N}\cdot\text{m}^2/\text{C}^2$.
Clamps $r$ to $10^{-10}$ near axis to avoid division by zero.

\subsubsection*{\texttt{EMethods} Enum}
\begin{lstlisting}[style=python]
class EMethods(Enum):
    FW  = (fw_e,)
    BOB = (bob_e,)

# Usage:
method = EMethods["FW"].value[0]   # -> fw_e callable
\end{lstlisting}

\subsection{\texttt{eext.py}}

\subsubsection*{\texttt{fill\_eext\_file(filepath, method, dia, offset, Q, L, N)} $\to$ \texttt{Path}}

Main public function. Returns immediately if renamed output already exists; otherwise:

\begin{enumerate}
    \item \texttt{get\_e\_field\_data(method, ...)} $\to$ $E_x$, $E_y$, $E_z$ shape $(N,N,N)$
    \item \texttt{\_fill\_efield\_datasets(...)} overwrites E datasets in HDF5
    \item \texttt{filepath.rename(new\_filepath)} appends E-field parameters to filename
\end{enumerate}

\subsubsection*{Per-point inner loop in \texttt{get\_e\_field\_data}}

For every grid point $(i,j,k)$ and each of the 6 coils:
\begin{lstlisting}[style=plain]
orient_point(coil, cartesian_point)
    -> subtract coil.position
    -> apply coil.orientation.inv()
toCyl(local_cartesian) -> (r, phi, z)
method(r, z, a, Q)     -> (E_r, E_z)   [E_phi = 0 by symmetry]
toCart(E_r, phi, E_z)  -> (dEx, dEy, dEz)
accumulate into Ex[i,j,k], Ey[i,j,k], Ez[i,j,k]
\end{lstlisting}

\noindent\textbf{Performance:} $\mathcal{O}(N^3 \times 6)$ non-vectorised calls.
At $N=72$: ${\approx}2.2\times10^6$ evaluations. Progress logged every 10\%.

% ============================================================
\section{Module: \texttt{src/utils}}

\subsection{\texttt{cyl.py}}

\subsubsection*{\texttt{toCyl(coord)} $\to$ \texttt{ndarray([$\rho$, $\varphi$, $z$])}}

\[
\rho = \sqrt{x^2+y^2},\quad
\varphi = \mathrm{arctan2}(y,x),\quad
z = z
\]

\subsubsection*{\texttt{toCart(r, theta, z)} $\to$ \texttt{(x, y, z)}}

\[
x = r\cos\theta,\quad y = r\sin\theta,\quad z = z
\]

\noindent\textbf{Note:} \texttt{theta} is the azimuthal angle of the \emph{position}
(from \texttt{toCyl}), not of the vector. Correct because $E_\varphi = 0$ by symmetry.

\subsection{\texttt{paths.py}}

\begin{lstlisting}[style=python]
_script_dir = Path(__file__).resolve().parent   # src/utils/
ROOT_DIR    = _script_dir.parent.parent          # WarpX/
BEXT_DIR    = ROOT_DIR / "output" / "bext"      # WarpX/output/bext/
\end{lstlisting}

\begin{longtable}{L{3cm} L{3.5cm} L{6.5cm}}
\toprule
\textbf{Constant} & \textbf{Resolves to} & \textbf{Used by} \\
\midrule
\texttt{ROOT\_DIR} & Project root     & Internal derivation \\
\texttt{BEXT\_DIR} & \texttt{output/bext/} & \texttt{src/bext/bext.py} \\
\bottomrule
\end{longtable}

\subsubsection*{Adding a new path}
\begin{lstlisting}[style=python]
# In paths.py
NEW_DIR = ROOT_DIR / "output" / "new_dir"

# Elsewhere
from src.utils.paths import NEW_DIR
\end{lstlisting}

% ============================================================
\section{Physics Background: Polywell Fusion}

\subsection{The Polywell Concept}

A polywell is an inertial electrostatic confinement (IEC) fusion device combining:
\begin{enumerate}
    \item \textbf{Magnetic cusps} from a polyhedral coil arrangement --- create a
          magnetic well that confines electrons near the centre.
    \item \textbf{Electrostatic potential well} --- the confined electron cloud builds
          negative space charge that electrostatically accelerates positive ions inward.
\end{enumerate}

\subsection{Magnetic Configuration}

Six coils at $\pm x$, $\pm y$, $\pm z$ carry alternating currents (Table~\ref{tab:coils}).
The resulting field forms magnetic mirrors at the cusps; electrons are reflected back by
the stronger field at the cusp regions.

\subsection{Electrostatic Field}

Modelled as six uniformly charged rings (same geometry as coils).
The field at cylindrical position $(r,z)$ from a ring of radius $a$ and charge $Q$:

\begin{align}
E_r &= \frac{\lambda a}{4\pi\varepsilon_0}
       \int_0^{2\pi} \frac{r - a\cos\theta}{D^3}\,d\theta \\[6pt]
E_z &= \frac{\lambda a\,z}{4\pi\varepsilon_0}
       \int_0^{2\pi} \frac{1}{D^3}\,d\theta \\[4pt]
\lambda &= \frac{Q}{2\pi a}, \qquad
D = \sqrt{r^2 + a^2 - 2ar\cos\theta + z^2}
\end{align}

\subsection{PIC Simulation Setup}

\begin{longtable}{L{5cm} L{9cm}}
\toprule
\textbf{Property} & \textbf{Value} \\
\midrule
Grid               & 3D Cartesian $[-L,L]^3$ \\
Field BCs          & Open \\
Particle BCs       & Absorbing \\
Solver             & Electromagnetic, Yee scheme, CFL = 0.99 \\
Species            & Electrons + Protons \\
Initial distribution & Uniform within $\pm 0.11L$ per axis \\
Initial velocity   & $0.9c$ in $z$-direction \\
External fields    & Pre-computed B and E, fixed (read-only) \\
Divergence cleaning & Disabled (\texttt{do\_initial\_div\_cleaning = 0}) \\
\bottomrule
\end{longtable}

\subsection{Key Physical Scales}

\begin{longtable}{L{4cm} L{4cm} L{6cm}}
\toprule
\textbf{Quantity} & \textbf{Value} & \textbf{Notes} \\
\midrule
Domain size       & $6\,\mathrm{m}^3$   & $2L$ per side \\
Coil radius       & 0.5 m              & \texttt{b\_dia} / 2 \\
Coil current      & 1 MA               & drives B-field \\
Plasma density    & $10^{18}\,\mathrm{m}^{-3}$ & \\
Time step         & 1 ps               & \texttt{const\_dt} \\
Total sim. time   & 10 ns              & \texttt{max\_steps} $\times$ \texttt{const\_dt} \\
\bottomrule
\end{longtable}

\subsection{References}

\begin{itemize}
    \item Bussard, R.W. (1991). ``Some Physics Considerations of Magnetic
          Inertial-Electrostatic Confinement.'' \textit{Fusion Technology}, 19(2).
    \item WarpX documentation: \url{https://ecp-warpx.github.io/}
    \item magpylib documentation: \url{https://magpylib.readthedocs.io/}
    \item openPMD standard: \url{https://github.com/openPMD/openPMD-standard}
\end{itemize}

\end{document}
